%% ========================================================================
%% Engineering Applications of Artificial Intelligence (EAAI)
%% Elsevier elsarticle class — preprint 3p (single-column review format)
%% Restructured: Intro → Related Work → Methodology → Experiments → Conclusion
%% ========================================================================
\documentclass[preprint,12pt,3p]{elsarticle}

\usepackage{lineno,hyperref}
\usepackage{xspace}
\usepackage{amsmath,amssymb,amsthm}
\usepackage{booktabs}
\usepackage{multirow}
\usepackage{graphicx}
\graphicspath{{figures/}{figures/paper/}}
\usepackage{xcolor}
\usepackage{algorithm}
\usepackage{algorithmic}
\usepackage{subfigure}
\usepackage{float}
\usepackage{microtype}
\usepackage{cleveref}
\usepackage{indentfirst}
\usepackage{bookmark}
%% EAAI margin control (A4, per Elsevier author guidelines).
%% geometry is already loaded by the class; use \geometry{} to reconfigure.
\geometry{a4paper, top=2.5cm, bottom=2.5cm, left=2.5cm, right=2.5cm}
\emergencystretch=2em
\hfuzz=1pt
\raggedbottom
\hyphenation{mag-net-o-stric-tion mag-net-o-stric-tion-in-duced
             mi-cro-con-trol-ler non-sta-tion-ar-i-ty}

%% Fix elsarticle \paragraph: restore \parindent and avoid double period
%% (elsarticle hardcodes 0pt indent and auto-appends "." to the heading)
\makeatletter
\renewcommand\elsparagraph{\@startsection{paragraph}{4}{\parindent}%
           {10\p@ \@plus 6\p@ \@minus 3\p@}%
           {-6\p@}%
           {\normalfont\itshape}}
\makeatother

\hypersetup{colorlinks, linkcolor=blue, anchorcolor=blue, citecolor=green}

\journal{Engineering Applications of Artificial Intelligence}
\bibliographystyle{elsarticle-harv}
\biboptions{authoryear}
\let\cite\citep   %% make \cite parenthetical: (Author, year)

%% ---- Custom macros ----
\newcommand{\zerone}{\textsc{SpecRas}\xspace}      %% our Spectral Raster method
\newcommand{\hetspatvae}{\textsc{HetSpatVAE}\xspace} %% our Heterogeneous Spatial VAE
\newcommand{\hczv}{\hetspatvae}                    %% backward compat alias
\newcommand{\R}{\mathbb{R}}
\newcommand{\E}{\mathbb{E}}
\newcommand{\bx}{\mathbf{x}}
\newcommand{\bs}{\mathbf{s}}
\newcommand{\bI}{\mathbf{I}}
\newcommand{\bz}{\mathbf{z}}
\newcommand{\bmu}{\boldsymbol{\mu}}
\newcommand{\bSigma}{\boldsymbol{\Sigma}}
\newcommand{\ie}{i.e.,\xspace}
\newcommand{\eg}{e.g.,\xspace}
\newcommand{\etal}{\textit{et al.}\xspace}

\begin{document}

\begin{frontmatter}

%% ---- Title ----
%% Author names, affiliations, and correspondence are in titlepage.tex.
\title{Dual-Pathway Deep Learning for Power Transformer Vibration Fault
       Diagnosis: Spectral Feature Encoding and Spatial Variational
       Autoencoder}

%% ---- Highlights ----
%% 3–5 bullet points, each ≤ 85 characters (EAAI requirement).
\begin{highlights}
\item Validated on 19 real in-service transformers with strict
      cross-unit test split (no unit in more than one split).
\item SpecRas spectral rasterisation with Grad-CAM yields
      Hz-resolution spectral attribution maps.
\item HetSpatVAE detects anomalies label-free via spatial
      reconstruction residuals over the time-frequency plane.
\item SWA adds $+2.75$\,pp test accuracy with zero missed faults
      under cross-unit covariate shift.
\item Peak-load false positives suggest load metadata as the
      primary missing covariate.
\end{highlights}

%% ---- Abstract ----
\begin{abstract}
\indent Vibration-based fault diagnosis for in-service power transformers faces
three compounding challenges: non-stationary harmonic spectra under
varying load, severe label scarcity (a handful of confirmed faults per
unit per decade), and the need for frequency-level interpretability
before engineers can justify maintenance decisions.

This paper presents a dual-pathway deep learning framework validated on
19 real in-service alternating current (AC) power transformers,
partitioned by unit identity (15 training / 2 validation / 2 held-out
test; 400 test segments: 246 normal, 154 fault).
The supervised \zerone{} (Spectral Rasterisation) pathway rasterises a
1{,}200-dimensional spectral feature vector---time-domain statistics,
Welch power spectral density (PSD) at 1--2{,}000\,Hz, and
high-frequency energy ratios---into a $150{\times}150$ image and
fine-tunes a ResNet18 classifier with stochastic weight averaging (SWA);
Gradient-weighted Class Activation Mapping (Grad-CAM) converts
predictions to per-Hz spectral attribution maps.
The \hczv{} (Heterogeneous Spatial Variational Autoencoder) pathway
encodes three-channel time-frequency images (Morlet continuous wavelet
transform (CWT), short-time Fourier transform (STFT), waveform context)
through a spatial ResNet variational autoencoder; reconstruction
residuals localise anomalies without fault labels.
The \textit{AI contributions} are spectral feature rasterisation,
a spatial VAE architecture, and SWA-based training; the
\textit{engineering application} is interpretable fault diagnosis on
in-service power transformers under cross-unit covariate shift.

On 400 held-out samples from two previously unseen transformers
(246 normal, 154 fault), \zerone achieves \textbf{97.00\% accuracy}
(macro-F1\,=\,0.969, AUC\,=\,0.989) with zero missed faults; \hczv
reaches \textbf{98.25\% accuracy} (macro-F1\,=\,0.982) with zero
missed faults and just 7 false positives.
All false positives cluster in the evening peak-load window with
elevated 1{,}000--2{,}000\,Hz STFT content, consistent with load as
a missing covariate.
\end{abstract}

\begin{keyword}
condition monitoring \sep deep learning \sep industrial AI \sep
transformer fault diagnosis \sep variational autoencoder \sep
vibration signal analysis
\end{keyword}

\end{frontmatter}

\bookmark[page=1,level=0]{Dual-Pathway Deep Learning for Power Transformer Vibration Fault Diagnosis}
%% ========================================================================
\section{Introduction}
\label{sec:intro}
%% ========================================================================

A failed AC power transformer can black out a region for
hours; replacement takes months to years because each unit is largely
custom-built \cite{cigre2015TB642}.
Knowing whether a unit is degrading---before it fails---is therefore
worth considerable effort.
Surface-mounted accelerometers have become the standard non-intrusive
monitoring modality: they can be retrofitted without de-energising the
unit and are sensitive to core looseness and winding deformation at a
fraction of the cost of dissolved gas analysis (DGA) oil sampling
\cite{bagheri2018transformer,tenbohlen2016diagnostic}.

Three practical constraints make vibration-based transformer diagnosis
harder than the bear\-ing-fault literature suggests.

\textit{Non-stationarity.}
Transformer vibration stems from two coupled mechanisms:
magnet\-o\-stric\-tion-in\-duced Lorentz forces at even harmonics of the supply
frequency (100, 200\,Hz for a 50\,Hz system), and load-induced winding
forces at the same harmonics plus a broadband floor
\cite{zollanvari2021transformer}.
Both shift continuously with operating conditions.
A fault redistributes energy across frequency decades---but so does a
routine load step, and the two look similar in a fixed-window Fourier
representation \cite{hong2021transformer}.

\textit{Label scarcity.}
A large regional grid may see only a handful of confirmed fault cases
per transformer per decade.
Large, balanced, annotated training sets are simply not available, and
accumulating matched fault recordings for every unit is impractical
\cite{zhao2020shrinkage}.

\textit{Interpretability.}
Grid engineers are accountable for every maintenance decision.
An anomaly score alone cannot justify taking a high-value asset offline;
the model must indicate which frequency bands drove the conclusion so
engineers can cross-check against known physical fault signatures
\cite{rudin2019stop}.

\textit{Limitations of existing approaches.}
Rule-based methods encode domain expertise as spectral thresholds and
are fully transparent, but their accuracy is bounded by calibration
conditions and they cannot adapt to unseen fault morphologies
\cite{ieee2019c57}.
Empirical mode decomposition (EMD) and wavelet decomposition track
non-stationarity better but require manual feature selection \cite{peng2004wavelet}.
Convolutional neural network (CNN) classifiers on spectrograms achieve
high accuracy on rotating-machinery benchmarks but require balanced
labelled data and produce no frequency-level explanation
\cite{guo2019transfer,wen2018cnn}.
Flat-vector variational autoencoder (VAE) anomaly detectors avoid label
requirements but cannot
localise which region of the time-frequency image is anomalous
\cite{kingma2013auto,an2015variational}.

\textit{Contributions.}
This paper makes four technical contributions (Fig.~\ref{fig:overview}).
\begin{enumerate}
\item The \zerone{} (Spectral Rasterisation) pathway rasterises a
1{,}200-dimensional spectral feature vector (time-domain statistics,
Welch power spectral density (PSD) at 1--2\,000\,Hz, and
high-frequency energy ratios) into a $150{\times}150$ image whose
PSD-dominant strip preserves frequency ordering.  Following the
feature-to-image encoding of Zerone~\cite{kim2025zerone} and
substituting physics-motivated spectral descriptors for transformer
vibration, Gradient-weighted Class Activation Mapping (Grad-CAM)
activations in the PSD rows read directly as per-Hz spectral
attribution maps.
\item The \hczv{} (Heterogeneous Spatial Variational Autoencoder) pathway
encodes a three-channel heterogeneous time-frequency image into a
$64{\times}7{\times}7$ latent map via a Spatial ResNet VAE
(SpatialResNetVAE); each spatial cell corresponds to a specific
time-frequency region, so reconstruction residuals serve as
interpretable anomaly heat maps.
\item Stochastic weight averaging \cite{izmailov2018averaging} contributes
$+2.75$\,pp test accuracy over the best-validation checkpoint on the
3{,}844-sample training pool under cross-unit covariate shift, with
zero missed faults.
\item A failure-mode analysis reveals that all false positives from both
pathways cluster in the 16:00--20:30 peak-load window and carry
elevated 1\,000--2\,000\,Hz short-time Fourier transform (STFT) content,
pointing to load metadata conditioning as a concrete next step.
\end{enumerate}

Section~\ref{sec:related} surveys related work.
Section~\ref{sec:methods} describes the methodology.
Section~\ref{sec:experiments} reports experiments and discussion.
Section~\ref{sec:conclusion} concludes.


%% ========================================================================
\section{Related Work}
\label{sec:related}
%% ========================================================================

\subsection{Vibration-Based Transformer Fault Diagnosis}

Transformer condition monitoring has traditionally centred on dissolved
gas analysis and power-frequency electrical tests
\cite{ieee2019c57,cigre2015TB642},
both requiring oil sampling or controlled energisation.
Vibration analysis emerged as a non-intrusive complement
\cite{bagheri2018transformer}: core looseness, winding
deformation, and cooling faults each perturb the 100--200\,Hz harmonic
components in recognisable ways.
Early systems extracted these features by hand and applied fixed
thresholds or classical classifiers
\cite{tenbohlen2016diagnostic,zollanvari2021transformer}.
More recent work uses deep recurrent networks on waveform windows and
CNN classifiers on vibration spectrograms
\cite{hong2021transformer,li2023cnn}, with
reported accuracies above 90\% on within-unit test sets.
Cross-unit evaluation---where test transformers are entirely absent
from training---is less commonly reported and consistently harder.

\subsection{Deep Learning for Industrial Condition Monitoring}

CNN classifiers on spectrograms are now standard for rotating-machinery
fault diagnosis, reaching near-perfect accuracy on benchmarks such as
Case Western Reserve University (CWRU) and Jiangnan University (JNU)
benchmarks \cite{guo2019transfer,wen2018cnn}.
Attention mechanisms and recurrent architectures extend the temporal
modelling horizon \cite{zhao2019deep,zhao2017bilstm}.
Transfer learning from ImageNet cuts the labelled-data requirement
substantially \cite{shao2019transfer}; deep residual
shrinkage networks add robustness to noisy labels
\cite{zhao2020shrinkage}.
Power transformer vibration differs from bearing diagnostics in two ways.
The training set is far smaller---hundreds of samples per class, not
thousands.
And the signal concentrates in narrow harmonic peaks at multiples of
the supply frequency, so raw-signal CNNs without explicit frequency-domain
feature engineering tend to miss the discriminative structure.

\subsection{Signal Representation and Feature Learning}

How a signal is represented before it reaches a model has a larger effect
on diagnostic accuracy than the model architecture itself, particularly
when training data are limited.
Continuous wavelet transform (CWT) resolves transient events at multiple
time-frequency scales and has been standard for transformer fault detection
\cite{peng2004wavelet}.
STFT is preferred for quasi-stationary harmonic signals because of its
computational efficiency and fixed spectral resolution
\cite{cohen1995tfa}.
Combining CWT and STFT captures complementary fault signatures
\cite{lei2020review}.
Signal-to-image encoding for CNN processing ranges from raw waveform
reshaping \cite{wen2018cnn} to stacked short-time spectra
as image channels \cite{zhao2019deep}.
End-to-end 1-D CNNs learn task-relevant frequency filters without
manual selection \cite{guo2019transfer}.
Our \zerone encoding adapts the feature-to-image rasterisation principle
of Zerone~\cite{kim2025zerone} to physics-motivated spectral
features specific to power transformer vibration.
Rather than mapping binary or generic indicators, \zerone encodes a
1{,}200-dimensional spectral feature vector (Welch PSD at
1--2\,000\,Hz, STFT frame means, and high-frequency energy ratios)
into a physically ordered raster image.
The PSD panel occupies $\approx$77\% of image height with frequency
ordering preserved row-by-row, so Grad-CAM activations read directly
as Hz-resolution spectral attribution maps.

\subsection{Generative Models for Anomaly Detection}

VAEs \cite{kingma2013auto} are well suited to one-class anomaly
detection: a model trained exclusively on normal samples produces high
reconstruction error and out-of-distribution latent codes for anomalous
inputs \cite{an2015variational}.
The $\beta$-VAE \cite{higgins2017betavae} promotes disentangled
representations at the cost of reconstruction fidelity.
Deep Autoencoding Gaussian Mixture Models \cite{zong2018deep} combine
reconstruction error with latent-space density estimation; PatchCore
\cite{roth2022towards} builds a memory bank of normal feature patches
for industrial image anomaly detection.
Both use flat vector latent codes and therefore cannot say where in the
time-frequency image the anomaly occurred.
\hczv addresses this by retaining the $7{\times}7$ spatial structure
of the ResNet encoder, with posterior collapse managed via $\beta$-KL
warm-up annealing \cite{lucas2019posterior}.


%% ========================================================================
\section{Methodology}
\label{sec:methods}
%% ========================================================================

%% --- Figure: dual-pathway framework overview ---
\begin{figure}[H]
  \centering
  \subfigure[\zerone supervised pathway: 1\,200-dim spectral feature vector
    $\to$ $150\times150$ raster image $\to$ ResNet18\,+\,SWA.
    Best-val checkpoint (V1) reaches 94.25\%; SWA checkpoint (V3) reaches 97.00\%.]{%
    \includegraphics[width=0.88\linewidth]{spec_feat.png}%
  }\\[8pt]
  \subfigure[\hczv one-class pathway: data preprocessing $\to$
    SpatialResNetVAE training $\to$ composite anomaly scoring.
    Trained on normal samples only; achieves 98.25\% accuracy.]{%
    \includegraphics[width=0.72\linewidth]{vae_flow.png}%
  }
  \caption{%
    Dual-pathway vibration fault diagnosis framework.
    \textbf{(a)} \zerone supervised pathway: spectral rasterisation
    $\to$ ResNet18\,+\,SWA $\to$ per-Hz Grad-CAM attribution.
    \textbf{(b)} \hczv one-class pathway: 3-channel time-frequency image
    $\to$ SpatialResNetVAE $\to$ reconstruction-residual heat map.
    Both branches share energy-weighted multi-channel fusion
    (Section~\ref{sec:preprocess}).
  }
  \label{fig:overview}
\end{figure}

\subsection{Problem Formulation}
\label{sec:problem}

Let $\mathcal{D} = \{(\mathcal{X}_k, y_k)\}_{k=1}^{N}$ be a labelled
dataset of vibration samples, where
%
\begin{equation}
  \mathcal{X}_k = \bigl\{
    x_k^{(u)} \in \R^{T} \;\big|\; u = 1, \ldots, U_k
  \bigr\}
  \label{eq:sample_def}
\end{equation}
%
is the set of multi-channel signals acquired at timestamp $t_k$,
$T = 8{,}192$ is the fixed signal length (sampling rate
$f_s = 8{,}192\,$Hz, 1-second window), $U_k \geq 1$ is the number of
active sensor channels, and $y_k \in \{0,1\}$ is the binary label
(0\,=\,normal, 1\,=\,fault).
The diagnostic task has two objectives: (i) supervised binary
classification given $\mathcal{X}_k$ and a labeled training set; and
(ii) one-class anomaly detection given only normal-class training data
via a composite scoring function.
Both require interpretable evidence linking the decision to specific
frequency bands.

\subsection{Signal Preprocessing}
\label{sec:preprocess}

\paragraph{Energy-weighted channel fusion}
To collapse $\mathcal{X}_k$ into a single representative signal,
the channel weights proportional to the signal energy are calculated:
%
\begin{align}
  w_u &= \frac{
    \E\!\left[(x_k^{(u)})^2\right]
  }{
    \displaystyle\sum_{j=1}^{U_k}
      \E\!\left[(x_k^{(j)})^2\right] + \varepsilon
  }, \label{eq:energy_weight} \\[4pt]
  \tilde{x}_k &= \sum_{u=1}^{U_k} w_u \, x_k^{(u)}. \notag
\end{align}
%
with $\varepsilon = 10^{-12}$.
This fusion is invariant to channel ordering and count, so all
downstream encoders remain agnostic to the sensor configuration of
each individual transformer.
The acquired JSON records contain only a \texttt{sensor\_id} string
and signal samples (no position, mounting angle, or sensor type is
recorded), so energy-weighted aggregation is the natural choice for
producing a fixed-size representation regardless of how many channels
happen to be active.

\paragraph{DC removal and label assignment}
Both pathways operate on the mean-subtracted signal
$\bar{x}_k = \tilde{x}_k - \mu_{\tilde{x}_k}$.
Class labels are inferred automatically from directory names:
subdirectories containing normal receive label 0; those
containing fault receive label 1.


\subsection{Supervised \zerone--ResNet18 Pathway}
\label{sec:zerone}

\subsubsection{Spectral Feature Vector}
\label{sec:feat_vec}

For each DC-removed channel signal $\bar{x}^{(u)} \in \R^T$, a
$d = 1{,}200$-dimensional characteristic vector is extracted following the
schema in Table~\ref{tab:feat_schema}.

\begin{table}[H]
  \centering
  \small
  \caption{%
    \zerone feature schema ($d = 1{,}200$ total dimensions).
    Multi-channel results are aggregated by energy-weighted mean
    (Eq.~\ref{eq:energy_weight}).
  }
  \label{tab:feat_schema}
  \renewcommand{\arraystretch}{1.2}
  \begin{tabular}{@{}llrp{8cm}@{}}
    \toprule
    Segment & Method & Dims & Description \\
    \midrule
    Time   & Statistics  & 15    &
      Mean, RMS, variance, standard deviation, maximum, minimum,
      peak-to-peak, kurtosis, skewness, zero-crossing rate,
      mean absolute value, crest factor, impulse factor,
      margin factor, and waveform factor \\
    STFT   & Frame means & 127   &
      Mean magnitude per frame
      ($n_\mathrm{seg}{=}128$, hop${=}64$, DC bin removed) \\
    PSD    & Welch       & 1\,050 &
      1--1\,000\,Hz at 1\,Hz resolution (1\,000 bins)
      $+$ 1\,001--2\,000\,Hz in 20\,Hz bands (50 bins) \\
    HF     & Ratios      & 8     &
      Amplitude ratio and power ratio above each of
      $f_\mathrm{thr} \in \{1000,2000,3000,4000\}$\,Hz \\
    \midrule
    \multicolumn{2}{@{}l}{\textbf{Total}} & \textbf{1\,200} & \\
    \bottomrule
  \end{tabular}
\end{table}

\paragraph{Time-domain features}
For signal $\bar{x} \in \R^T$, the 15 time-domain statistics are
\begin{equation}
\begin{aligned}
\mathbf{z}_{\mathrm{time}}
=
\bigl[
&\mu,\,
x_\mathrm{RMS},\,
\sigma^2,\,
\sigma,\,
x_{\max},\\
&x_{\min},\,
x_\mathrm{p2p},\,
\kappa,\,
\gamma,\,
r_\mathrm{zc},\\
&\mu_{|\cdot|},\,
x_\mathrm{crest},\,
x_\mathrm{imp},\\
&x_\mathrm{margin},\,
x_\mathrm{wave}
\bigr]
\in \R^{15}.
\end{aligned}
\end{equation}
Let
\begin{align}
\mu &= T^{-1}\sum_{t=1}^{T}\bar{x}_t, \notag \\
\sigma^2 &= T^{-1}\sum_{t=1}^{T}(\bar{x}_t-\mu)^2, \notag \\
\sigma &= \sqrt{\sigma^2}.
\end{align}
Then the main quantities are defined as
\begin{equation}
\begin{aligned}
x_\mathrm{RMS} \;&=\; \sqrt{T^{-1}\sum_{t=1}^{T}\bar{x}_t^2},\\
x_{\max} \;&=\; \max_t \bar{x}_t,\qquad
x_{\min} \;=\; \min_t \bar{x}_t,\\
x_\mathrm{p2p} \;&=\; x_{\max}-x_{\min},\\
\kappa \;&=\; \frac{m_4}{m_2^2},\\
\gamma \;&=\; \frac{m_3}{\sigma^3},\\
r_\mathrm{zc} \;&=\; \frac{1}{T-1}
  \sum_{t=1}^{T-1}\mathbf{1}\!\left[\bar{x}_t\,\bar{x}_{t+1}<0\right],\\
\mu_{|\cdot|} \;&=\; T^{-1}\sum_{t=1}^{T}|\bar{x}_t|,\\
x_\mathrm{crest} \;&=\; \frac{x_{\max}}{x_\mathrm{RMS}},\\
x_\mathrm{imp} \;&=\; \frac{x_{\max}}{\mu_{|\cdot|}},\\[2pt]
x_\mathrm{wave} \;&=\; \frac{x_\mathrm{RMS}}{\mu_{|\cdot|}},\\
x_\mathrm{margin} \;&=\; \frac{x_{\max}}
{\left(T^{-1}\sum_{t=1}^{T}|\bar{x}_t|^4\right)^{1/4}}.
\end{aligned}
\label{eq:time_features}
\end{equation}
where $m_p = T^{-1}\sum_{t=1}^{T}(\bar{x}_t-\mu)^p$ denotes the
$p$-th central moment.

\paragraph{STFT features}
After DC removal, a short-time Fourier transform is applied with
Hann window, segment length $n_\mathrm{seg}=128$, and overlap
$n_\mathrm{overlap}=64$. Let $S(f,\tau)$ denote the complex STFT of
$\bar{x}$. After discarding the DC bin, we compute the mean magnitude
of each time frame as
\begin{equation}
  s_\tau = \frac{1}{|\mathcal{F}^\prime|}
  \sum_{f \in \mathcal{F}^\prime} |S(f,\tau)|,
\end{equation}
where $\mathcal{F}^\prime$ is the non-DC frequency index set. The
resulting sequence is zero-padded or truncated to 127 dimensions,
yielding a fixed-length STFT descriptor
\begin{equation}
\mathbf{s}_{\mathrm{STFT}} \in \R^{127}.
\end{equation}

\paragraph{PSD features}
Welch PSD is estimated at $1$\,Hz resolution from $1$ to $2000$\,Hz.
The low-frequency region $1$--$1000$\,Hz is retained at full
resolution, while the high-frequency region $1001$--$2000$\,Hz is
compressed into 20\,Hz bands by local averaging. This gives
\begin{equation}
\begin{aligned}
\mathbf{p}_{\mathrm{PSD}}
=
\bigl[
&P(1),\,P(2),\,\dots,\,P(1000),\\
&\bar{P}_{1001{:}1020},\,\dots,\,\bar{P}_{1981{:}2000}
\bigr]
\in \R^{1050}.
\label{eq:welch}
\end{aligned}
\end{equation}

\paragraph{High-frequency ratio features}
To explicitly characterise the proportion of elevated-frequency
content, we define four threshold frequencies
\[
f_\mathrm{thr} \in \{1000, 2000, 3000, 4000\}\,\text{Hz}.
\]
For each threshold, two scalar ratios are computed: the amplitude
ratio
\begin{equation}
  r^{(A)}(f_\mathrm{thr})
  =
  \frac{\sum_{f \ge f_\mathrm{thr}} A(f)}
       {\sum_{f \ge 0} A(f)},
\end{equation}
and the power ratio
\begin{equation}
  r^{(P)}(f_\mathrm{thr})
  =
  \frac{\sum_{f \ge f_\mathrm{thr}} P(f)}
       {\sum_{f \ge 0} P(f)},
\end{equation}
where $A(f)$ and $P(f)$ denote the spectrum magnitude and PSD,
respectively. Concatenating the four threshold pairs yields
\begin{equation}
\begin{aligned}
\mathbf{r}_{\mathrm{HF}}
=
\bigl[
&r^{(A)}(1000),\,r^{(P)}(1000),\\
&r^{(A)}(2000),\,r^{(P)}(2000),\\
&r^{(A)}(3000),\,r^{(P)}(3000),\\
&r^{(A)}(4000),\,r^{(P)}(4000)
\bigr]
\in \R^{8}.
\end{aligned}
\end{equation}

\paragraph{Multi-channel aggregation}
For a sample with $U$ available sensor channels, each channel yields a
feature vector $\mathbf{z}_u \in \R^{1200}$. To preserve the relative
importance of stronger vibration responses, the sample-level feature is
computed by an energy-weighted mean:
\begin{equation}
\mathbf{z}
=
\frac{\sum_{u=1}^{U} e_u \mathbf{z}_u}
     {\sum_{u=1}^{U} e_u + \varepsilon},
\label{eq:aggregation}
\end{equation}
where $e_u = T^{-1}\sum_{t=1}^{T} x_{u,t}^2$ is the channel energy and
$\varepsilon$ is a small constant for numerical stability.
\subsubsection{Raster Stripe Image Encoding}
\label{sec:raster}

The aggregate feature vector $\mathbf{S}_k$ is mapped to a
$150 \times 150$ RGB image using a non-interpolating raster layout
(Fig.~\ref{fig:pipeline}).
Features are tiled into a 2D grid with unit cell size $W_\mathrm{unit}
\times H_\mathrm{unit} = 3 \times 2$ pixels, and the four panels are
stacked vertically in order (time / STFT / PSD / HF), with each panel
wrapping its features row-by-row at a fixed tile width.
%
\begin{equation}
  n_\mathrm{rows}^{(g)} = \left\lceil d_g / w_g \right\rceil,
  \label{eq:raster_rows}
\end{equation}
where $g \in \{\text{time},\text{stft},\text{psd},\text{hf}\}$,
$d_g$ is the segment dimension, and $w_g$ is the tile width
(16 / 16 / 32 / 8 for the four panels).
The PSD panel spans 33 rows $\times$ 2\,px $= 66$\,px,
against 2\,px for time and HF and 16\,px for STFT, so it
occupies approximately 77\% of the image content area.

Pixel values are linearly normalised using training-set statistics:
%
\begin{equation}
  I_{ij} = \mathrm{clip}\!\left(
    \frac{S_d - \hat{m}_d}{\hat{M}_d - \hat{m}_d + \varepsilon},\; 0, 1
  \right),
  \label{eq:norm}
\end{equation}
%
where $\hat{m}_d$ and $\hat{M}_d$ are training-set minimum and maximum
for dimension $d$.
The image is padded to $150{\times}150$ and represented as a
three-channel input tensor for ResNet18-based classification.

Two properties of this layout matter for interpretability and learning.
Each PSD row holds 32 tiles spanning 32\,Hz: row $r$ covers
$1 + 32r$ to $32(r+1)$\,Hz, so Grad-CAM activations in those rows read
directly as frequency bands.
The 77\% height share also concentrates most CNN receptive fields on
spectral content; the time and HF panels occupy narrow border strips.

\subsubsection{ResNet18 Training and SWA}
\label{sec:resnet}

A ResNet18 \cite{he2016deep} classifier is trained from scratch on
\zerone images.
The final fully-connected layer is replaced with a two-output linear
head ($\R^{512} \to \R^2$).
Table~\ref{tab:hyper_zerone} lists the complete hyperparameters.

\begin{table}[H]
  \centering
  \small
  \caption{%
    Training configuration for \zerone ResNet18.
    All values were fixed before any test-set evaluation.
  }
  \label{tab:hyper_zerone}
  \renewcommand{\arraystretch}{1.2}
  \begin{tabular}{@{}lp{3.8cm}@{}}
    \toprule
    Hyperparameter & Value \\
    \midrule
    Image size               & $150 \times 150$ \\
    Batch size               & 16 \\
    Max epochs               & 30 \\
    Base learning rate       & $8\times10^{-6}$ \\
    LR schedule              & ReduceLROnPlateau (val macro-F1) \\
    Weight decay             & $10^{-4}$ \\
    Label smoothing          & $\varepsilon = 0.05$ \\
    Class weights            & Inverse class frequency \\
    Gradient clipping        & $\|\mathbf{g}\|_2 \leq 1.0$ \\
    EMA decay                & $0.999$ \\
    SWA start ratio          & 0.50 of max epochs \\
    SWA min-LR               & $5 \times 10^{-5}$ \\
    SWA max-loss bias        & 0.05 above best val loss \\
    Random seed              & 42 \\
    \bottomrule
  \end{tabular}
\end{table}

\paragraph{Loss function}
Cross-entropy with label smoothing and class weights:
%
\begin{equation}
  \mathcal{L}_\mathrm{CE} =
  -\sum_{c=0}^{1} \alpha_c \,
  \bigl[(1-\varepsilon)\,y_c + \tfrac{\varepsilon}{2}\bigr]
  \log \hat{p}_c,
  \label{eq:ce_loss}
\end{equation}
%
where $\alpha_c \propto n_c^{-1}$ are inverse-frequency class weights.

\paragraph{Model selection}
The best checkpoint is selected by the four-key tie-break
\[
  \bigl(F1_\mathrm{val}\!\uparrow,\; L_\mathrm{val}\!\downarrow,\;
  \mathrm{lr}\!\downarrow,\; \mathrm{epoch}\!\uparrow\bigr).
\]
An EMA copy (decay 0.999) of the parameters runs in parallel.

\paragraph{Stochastic weight averaging}
In the second half of training, whenever the current validation loss
falls within 5\% of the best observed and the learning rate drops
below a threshold, the weights are accumulated into an SWA buffer
\cite{izmailov2018averaging}.
After training, batch-normalisation running statistics are
re-computed on 128 forward passes.

\paragraph{Grad-CAM attribution}
Grad-CAM \cite{selvaraju2017grad} is applied to the final residual
block (\texttt{layer4}).
The resulting $5{\times}5$ activation map is upsampled to
$150{\times}150$ and overlaid on the raster image.
Hot regions within the PSD panel (rows 1--33 of that panel) correspond
to specific frequency bands: a pixel at in-panel row $r$ (0-based),
column $c$ maps to $f = 1 + 32r + c$\,Hz, giving per-Hz spectral
attribution without additional decoding.

%% --- Figure: SpecRas feature extraction and raster-stripe pipeline ---
\begin{figure}[H]
  \centering
  \subfigure[Feature construction: 4 descriptor groups
    (time\,15 + STFT\,127 + PSD\,1\,050 + HF\,8 = 1\,200 dims)
    $\to$ $150\times150$ raster image.
    PSD occupies $\approx$77\% of the image; each PSD row spans 32\,Hz
    so Grad-CAM activations read directly as Hz bands.]{%
    \includegraphics[width=0.88\linewidth]{spec_feat.png}%
  }\\[8pt]
  \subfigure[ResNet18 backbone: four residual layers
    ($64\!\to\!128\!\to\!256\!\to\!512$ channels),
    global average pooling, two-class linear head
    ($\R^{512}\!\to\!\R^{2}$).]{%
    \includegraphics[width=0.82\linewidth]{resnet18.png}%
  }
  \caption{%
    \zerone pipeline.
    \textbf{(a)} Feature extraction and raster encoding:
    PSD dominates 77\% of the $150\times150$ image; Grad-CAM hot rows
    map directly to 1--1\,000\,Hz spectral bands.
    \textbf{(b)} ResNet18 classifier backbone.
  }
  \label{fig:pipeline}
\end{figure}

\subsection{\hczv Anomaly Detection Pathway}
\label{sec:vae}

\subsubsection{Heterogeneous Three-Channel Encoding}
\label{sec:3ch}

The energy-weighted z-score normalised signal
$\check{x} = (\tilde{x} - \mu) / \sigma$
is converted to a three-channel $224{\times}224$ image
$\bI \in [0,1]^{3 \times 224 \times 224}$.
Let
%
\begin{align}
  \mathbf{M}_0 &= \bigl|\mathrm{CWT}(\check{x},\,
                  \Psi_\mathrm{morl},\,a_{1:128})\bigr|,
                  \label{eq:cwt_ch} \\
  \mathbf{M}_1 &= \bigl|\mathrm{STFT}(\check{x};\,
                  n{=}256,\,h{=}128)\bigr|,
                  \label{eq:stft_ch} \\
  \mathbf{M}_2 &= \mathrm{fold}(\check{x},\,32{\times}256),
                  \label{eq:context_ch}
\end{align}
%
where scales $a_{1:128}$ are logarithmically spaced and
$\mathrm{fold}(\cdot)$ reshapes the 1-D signal into a $32{\times}256$
matrix.
Each channel is then:
%
\begin{equation}
  \bI_{[n]} = \mathrm{norm}\!\bigl(
    \mathrm{resize}(\mathbf{M}_n,\;224{\times}224)
  \bigr), \quad n=0,1,2.
  \label{eq:img_ch}
\end{equation}
The three channels provide complementary physical views: CWT resolves
transient signatures at multiple time-frequency scales; STFT reveals
stationary harmonic structure; context encodes global waveform
morphology.

\subsubsection{SpatialResNetVAE Architecture}
\label{sec:vae_arch}

\paragraph{Encoder}
A ResNet18 backbone processes $\bI$ through its stem and four residual
blocks, producing a feature tensor $\mathbf{h} \in \R^{512 \times 7 \times 7}$.
Two parallel $1{\times}1$ convolutional heads project $\mathbf{h}$ to
spatial mean and log-variance, both in $\R^{64 \times 7 \times 7}$:
%
\begin{align}
  \bmu_z &= \mathbf{W}_\mu * \mathbf{h}, \notag \\
  \log\boldsymbol{\sigma}_z^2 &= \mathbf{W}_\sigma * \mathbf{h}.
  \label{eq:spatial_proj}
\end{align}
%
Retaining the $7{\times}7$ spatial dimensions (rather than global
average pooling) is the single most important architectural decision:
each spatial cell $(j,k)$ has a $32{\times}32$-pixel receptive field
in $\bI$, so reconstruction residuals at cell $(j,k)$ correspond to
a specific time-frequency region of the input CWT or STFT.

\paragraph{Reparameterisation}
$\bz = \bmu_z + \boldsymbol{\sigma}_z \odot \boldsymbol{\epsilon}$,
$\boldsymbol{\epsilon} \sim \mathcal{N}(\mathbf{0},\mathbf{I})$.

\paragraph{Decoder}
A $1{\times}1$ convolution maps $\bz \in \R^{64\times 7\times 7}$
to 512 channels; five transposed-convolution blocks
(ConvTranspose2d, $k{=}4$, $s{=}2$, $p{=}1$, BatchNorm, ReLU)
progressively upsample to $3{\times}224{\times}224$.
The output activation is Sigmoid.

\paragraph{Training objective}
The $\beta$-VAE evidence lower bound (ELBO) \cite{kingma2013auto,higgins2017betavae}
with L1 reconstruction:
%
\begin{align}
  \mathcal{L} &=
  \frac{1}{B}\,\|\bI - \hat{\bI}\|_1\notag\\
  &\;+\;\beta(e)\cdot\frac{-1}{2B}
  \sum_{i,j,k}\bigl(
    1 + \log\sigma_{z,ijk}^2
    - \mu_{z,ijk}^2
    - \sigma_{z,ijk}^2
  \bigr),
  \label{eq:vae_loss}
\end{align}
%
where $\beta(e)=\min(\beta_{\max}, \beta_{\max}\cdot e/e_{\mathrm{warm}})$ anneals
from 0 to $\beta_{\max}=0.05$ over $e_{\mathrm{warm}}=20$ epochs.
L1 reconstruction preserves edge sharpness in the CWT scalogram channel;
the Kullback--Leibler (KL) warm-up prevents posterior collapse \cite{lucas2019posterior}.
Training uses only normal-labelled samples.

\subsubsection{Composite Anomaly Score}
\label{sec:anomaly}

\paragraph{Channel-weighted reconstruction error}
%
\begin{align}
  e_\mathrm{rec}(\bI) =\;&
  0.4\,\|\bI_{[0]} - \hat{\bI}_{[0]}\|_1 \notag \\
  &+ 0.5\,\|\bI_{[1]} - \hat{\bI}_{[1]}\|_1 \notag \\
  &+ 0.1\,\|\bI_{[2]} - \hat{\bI}_{[2]}\|_1.
  \label{eq:rec_err}
\end{align}
%
The STFT channel receives the highest weight (0.5) because structured
harmonic patterns are most sensitive to fault-induced changes.

\paragraph{Mahalanobis distance}
The spatial mean of the encoder output
$\bar{\bmu}_z = \frac{1}{49}\sum_{j,k} \bmu_{z,\cdot jk} \in \R^{64}$
provides a fixed-length embedding.
A Gaussian density is fitted to all training-set embeddings:
%
\begin{align}
  \bmu_\mathrm{tr} &= N^{-1}\textstyle\sum_n \bar{\bmu}_{z,n}, \\
  \bSigma_\mathrm{tr} &=
    N^{-1}\textstyle\sum_n
    (\bar{\bmu}_{z,n} - \bmu_\mathrm{tr})
    (\bar{\bmu}_{z,n} - \bmu_\mathrm{tr})^\top
    + \varepsilon\mathbf{I}, \\
  d_\mathrm{MD}(\bI) &= \sqrt{
    (\bar{\bmu}_z - \bmu_\mathrm{tr})^\top
    \bSigma_\mathrm{tr}^{-1}
    (\bar{\bmu}_z - \bmu_\mathrm{tr})
  }.
  \label{eq:md}
\end{align}

\paragraph{Composite score}
Both scores are z-normalised using training statistics, then combined:
%
\begin{align}
  a(\bI) &=
  0.6 \cdot \frac{e_\mathrm{rec} - \mu_\mathrm{rec}}
                  {\sigma_\mathrm{rec}} \notag \\
  &\quad {}+ 0.4 \cdot
  \frac{d_\mathrm{MD} - \mu_\mathrm{MD}}{\sigma_\mathrm{MD}}.
  \label{eq:score}
\end{align}
%
A sample is flagged when $a(\bI) \geq \tau$, with $\tau$ at the
97.5th percentile of training normal scores.
Table~\ref{tab:hyper_vae} lists all \hczv hyperparameters.

\begin{table}[H]
  \centering
  \small
  \caption{\hczv training hyperparameters.}
  \label{tab:hyper_vae}
  \renewcommand{\arraystretch}{1.2}
  \begin{tabular}{@{}lp{3.8cm}@{}}
    \toprule
    Hyperparameter & Value \\
    \midrule
    Image size                  & $224 \times 224$ \\
    Batch size                  & 32 \\
    Max epochs                  & 100 \\
    Learning rate               & $10^{-4}$ (Adam) \\
    Latent channels $C_z$       & 64 \\
    $\beta_{\max}$              & 0.05 \\
    KL warm-up epochs           & 20 \\
    Initialisation              & Random (no pre-training) \\
    Anomaly threshold quantile  & 97.5\% \\
    Reconstruction weight       & 0.6 \\
    Random seed                 & 42 \\
    \bottomrule
  \end{tabular}
\end{table}

\subsection{Method Comparison}
\label{sec:compare}

Table~\ref{tab:method_compare} places the two diagnostic approaches
side by side on the four axes that matter in practice: input
representation, training label requirement, model type, and
interpretability output.

\begin{table}[H]
  \centering
  \caption{%
    Side-by-side comparison of diagnostic approaches.
    \checkmark: required; --: not required.
  }
  \label{tab:method_compare}
  \small
  \renewcommand{\arraystretch}{1.3}
  \begin{tabular}{@{}p{2.9cm}p{2.0cm}p{3.0cm}p{2.3cm}cp{3.2cm}@{}}
    \toprule
    Approach
      & Paradigm
      & Input representation
      & Model
      & Labels?
      & Interpretability output \\
    \midrule

    \zerone{} + ResNet18 (SWA)
      & Supervised
      & 1\,200-dim PSD raster image
      & ResNet18 + SWA
      & \checkmark
      & Grad-CAM per-Hz heat map \\
    \hczv{} (SpatialResNetVAE)
      & One-class detection
      & 3-ch CWT\slash{}STFT\slash{}context image
      & $\beta$-VAE, spatial latent
      & --
      & Reconstruction residual heat map \\
    \bottomrule
  \end{tabular}
\end{table}

\zerone preserves the full spectral shape in a 1\,200-bin PSD-dominant
image at the cost of a labelled training set and sensitivity to
cross-unit covariate shift.
\hczv replaces the discriminative objective with a generative one:
learn to reconstruct normal time-frequency images, then flag anything
that reconstructs poorly.
No fault labels are required, and the spatial latent map localises the
anomaly to specific CWT/STFT cells.


%% ========================================================================
\section{Experiments and Discussion}
\label{sec:experiments}
%% ========================================================================

\subsection{Dataset}
\label{sec:dataset}

Vibration measurements were collected from real in-service AC power
transformers at a regional substation.
Each monitoring session is stored as a JSON file in which all records
sharing a common timestamp $t_k$ constitute one Segment.
Each record carries a \texttt{sensor\_id} string and a fixed-length
\texttt{signal\_value} array of $T=8{,}192$ samples at
$f_s=8{,}192\,$Hz (one second per Segment).
No sensor position, orientation, or type metadata is stored in the
collected data; the number of active channels $U_k$ per Segment
varies across sessions and transformer units.
Energy-weighted channel fusion (Section~\ref{sec:preprocess})
therefore collapses whatever channels are present into a single
representative signal, requiring no knowledge of sensor placement or
unit type.
Table~\ref{tab:dataset} summarises the train/validation/test split.
Fig.~\ref{fig:signal_chars} illustrates representative multi-channel
signal characteristics at an in-service substation, motivating the
energy-weighted fusion and the high-dimensional spectral feature schema.

\begin{table}[H]
  \centering
  \small
  \caption{%
    Dataset composition across 19 transformer units
    (4{,}499 Segments total).
    Each unit is assigned to exactly one split; no unit crosses
    split boundaries.
    Test units 134 and 135 were fully withheld from all
    development phases.
  }
  \label{tab:dataset}
  \renewcommand{\arraystretch}{1.2}
  \begin{tabular}{@{}lrrrp{4.8cm}@{}}
    \toprule
    Split      & Normal & Fault & Total & Transformer IDs \\
    \midrule
    Train      & 1{,}595 & 2{,}249 & 3{,}844 &
      120, 127--129, 143, 152 (normal);
      114, 132, 136, 153--155, 157--159 (fault) \\
    Validation & 91  & 164 & 255 & 146 (normal); 145 (fault) \\
    Test       & 246 & 154 & 400 & \textbf{134 (normal); 135 (fault)} \\
    \bottomrule
  \end{tabular}
\end{table}


%% --- Figure: multi-channel vibration signal characteristics ---
\begin{figure}[H]
  \centering
  \subfigure[Top-15 WPD frequency-band energy per sensor (7-min avg).
    Energy concentrates below 10\,Hz; inter-channel levels vary
    substantially across all six sensors.]{%
    \includegraphics[width=0.92\linewidth]{sig_wpd.png}%
  }\\[4pt]
  \subfigure[Velocity RMS vs.\ time with ISO\,10816 reference lines.
    Channels\,1 and\,3 exceed the upper caution threshold;
    channels\,4--6 lie in the normal band.]{%
    \includegraphics[width=0.46\linewidth]{sig_rms.png}%
  }\hfill
  \subfigure[Crest factor (peak/RMS) per channel.
    Channels\,5--6 show markedly elevated impulsive content ($>$3.2)
    versus 2.65--2.75 for channels\,1--4.]{%
    \includegraphics[width=0.46\linewidth]{sig_crest.png}%
  }\\[4pt]
  \subfigure[Per-channel normalised WPD heatmap (7-min avg).
    All channels concentrate energy in the 0--10\,Hz sub-band;
    channels\,5--6 show secondary excitation at 10--20\,Hz.]{%
    \includegraphics[width=0.92\linewidth]{sig_heat.png}%
  }
  \caption{%
    Multi-channel signal characteristics (representative in-service substation).
    \textbf{(a)} WPD band energy: energy concentrated below 10\,Hz,
    heterogeneous across 6 channels.
    \textbf{(b,c)} Velocity RMS trend and crest-factor distribution;
    channels\,5--6 breach caution thresholds.
    \textbf{(d)} Normalised WPD heatmap confirming low-frequency dominance.
    Inter-channel heterogeneity motivates energy-weighted fusion
    (Eq.~\ref{eq:energy_weight}).
  }
  \label{fig:signal_chars}
\end{figure}

All 19 transformer units are assigned to exactly one split: 15 to
training (6 normal-condition units and 9 fault-condition units), 2 to validation
(units 146 and 145), and 2 to the held-out test set (units 134
and 135).
This hard unit-level partition means that test accuracy measures
whether the learned spectral representation transfers to transformers
the model has never seen, not whether it interpolates within units
already in the training pool.
Test units 134 and 135 were not touched during feature design,
hyperparameter search, or threshold calibration, making them a
genuine out-of-sample evaluation.

\subsection{Implementation Details}
\label{sec:impl}

All experiments ran on a single workstation with an NVIDIA GPU
(CUDA 12.8), Python 3.11, PyTorch~2.7\allowbreak+cu128, and torchvision~0.22.
Scientific computing used NumPy 2.3, SciPy 1.16, and scikit-learn 1.5.
The random seed was fixed at 42.

\zerone training uses the hyperparameters in Table~\ref{tab:hyper_zerone};
training-set normalisation statistics for image encoding are computed
once from the training split and frozen.
\hczv training uses only normal-labelled training samples; the
Mahalanobis statistics and anomaly threshold are computed from
training-normal embeddings after the VAE converges.

\subsection{Comparison Methods}
\label{sec:baselines}

Three diagnostic configurations are evaluated (two pathways, with two checkpoints for \zerone):

\begin{enumerate}
  \item \textbf{\zerone + ResNet18 (best-val checkpoint):} the supervised
    pathway using the checkpoint with highest validation F1.
  \item \textbf{\zerone + ResNet18 (SWA checkpoint):} the same pipeline
    with the stochastic weight-averaged model.
  \item \textbf{\hczv (SpatialResNetVAE):} the one-class anomaly
    detection pathway trained without fault labels, evaluated using
    the 97.5th-percentile anomaly threshold.
\end{enumerate}

\subsection{Performance Comparison}
\label{sec:results}

Table~\ref{tab:main_results} reports five diagnostic metrics on the
400-sample held-out test set.
Fig.~\ref{fig:vae_results} shows the \hczv visual diagnostics.

\begin{table}[H]
  \centering
  \caption{%
    Diagnostic performance on the held-out test set
    (246 normal / 154 fault, $N=400$; transformers 134 and 135).
    Pre.(N): normal-class precision.
    Rec.(N/F): normal/fault recall.
    Best result per metric shown in \textbf{bold}.
    $^\dagger$: trained without fault labels.
  }
  \label{tab:main_results}
  \small
  \renewcommand{\arraystretch}{1.25}
  \begin{tabular}{@{}p{3.6cm}p{1.9cm}ccccc@{}}
    \toprule
    Method & Paradigm
      & Acc.$\uparrow$
      & Pre.(N)$\uparrow$
      & Rec.(N)$\uparrow$
      & Rec.(F)$\uparrow$
      & Macro-F1$\uparrow$ \\
    \midrule

    \zerone{} + ResNet18 (best-val ckpt)
      & Supervised   & 94.25\% & 1.000 & 0.907 & \textbf{1.000} & 0.941 \\
    \zerone{} + ResNet18 (SWA)
      & Supervised   & 97.00\% & \textbf{1.000}
                    & 0.951   & \textbf{1.000} & 0.969 \\
    \textbf{\hczv{} (SpatialResNetVAE)}$^\dagger$
      & One-class    & \textbf{98.25\%} & \textbf{1.000} & \textbf{0.972} & \textbf{1.000} & \textbf{0.982} \\
    \bottomrule
  \end{tabular}
\end{table}

\paragraph{Fault recall}
Both \zerone checkpoints and \hczv achieve fault recall of $1.000$:
all 154 fault samples are correctly flagged.
Zero false negatives across both deep-learning pathways is the
operationally important result: a missed fault can lead to thermal
runaway and equipment loss.

\paragraph{SWA narrows the false-positive gap}
The SWA checkpoint reduces the false-positive count from
$23/246 = 9.3\%$ to $12/246 = 4.9\%$ (normal recall $0.907 \to 0.951$).
Each false positive triggers a field inspection that costs operational
time and disrupts scheduled maintenance.
The result is consistent with the flat-minima hypothesis
\cite{izmailov2018averaging}: a wider loss basin generalises better
across the covariate shift between training and test transformer units.

\paragraph{\hczv performance}
Trained exclusively on normal samples, \hczv reaches 98.25\% accuracy
and macro-F1 of 0.982, surpassing the SWA supervised result by
1.25\,pp, with 7 false positives (2.8\% of normal).
The t-SNE projection in Fig.~\ref{fig:vae_results}(a) shows that test
fault samples are systematically displaced from the training-normal
cluster, confirming that the spatial VAE has learned a coherent
normality manifold.
Score histograms in Fig.~\ref{fig:vae_results}(c,d) reveal the
two well-separated score distributions: fault scores peak near 0.7
and normal scores above 3.5, with the threshold at $\tau=2.00$
cleanly between the two modes.
The confusion matrix (Fig.~\ref{fig:vae_results}(b)) confirms
239 TN, 7 FP, 0 FN, 154 TP.

%% --- Figure: HetSpatVAE diagnostic outputs ---
\begin{figure}[H]
  \centering
  \subfigure[Latent space (t-SNE): training normal (grey),
    test normal (blue, $\circ$), test fault (red, $\times$).
    Fault samples displaced from the normal cluster confirm
    a coherent normality manifold.]{%
    \includegraphics[width=0.46\linewidth]{vae_tsne.png}%
  }\hfill
  \subfigure[Confusion matrix: 239\,TN, 7\,FP, 0\,FN, 154\,TP.
    Accuracy\,=\,98.25\%; fault recall\,=\,1.000 (zero missed faults).]{%
    \includegraphics[width=0.46\linewidth]{vae_cm.png}%
  }\\[4pt]
  \subfigure[Score distribution, test-normal ($n=246$,
    $\tau=2.00$). Mode near 4.0; 7 samples fall below $\tau$
    (false-positive rate 2.8\%).]{%
    \includegraphics[width=0.46\linewidth]{vae_hn.png}%
  }\hfill
  \subfigure[Score distribution, test-fault ($n=154$, $\tau=2.00$).
    All fault scores lie below $\tau$ (mode $\approx$0.7):
    100\% fault recall.]{%
    \includegraphics[width=0.46\linewidth]{vae_hf.png}%
  }
  \caption{%
    \hczv diagnostic outputs on the held-out test set
    ($\tau$ at 97.5th training-normal percentile).
    \textbf{(a)} t-SNE: fault latent codes occupy a distinct region.
    \textbf{(b)} Confusion matrix: 98.25\% accuracy, zero missed faults.
    \textbf{(c,d)} Score histograms: fault and normal distributions
    are well separated at $\tau=2.00$.
  }
  \label{fig:vae_results}
\end{figure}

\paragraph{ROC, precision--recall, and calibration (\zerone SWA)}
Fig.~\ref{fig:zerone_results} shows the full diagnostic outputs for
the \zerone SWA checkpoint.
The ROC AUC reaches $0.989$ for both classes; the fault-class
average precision (AP) is $0.960$, so the ranked probability outputs
are reliable regardless of which threshold is chosen.
The reliability (calibration) curve gives an expected calibration
error (ECE) of $0.128$: a model output of 90\% fault probability
corresponds to an observed fault rate close to 90\%.
The t-SNE projection of penultimate-layer features shows two
well-separated clusters, and the per-band PSD mean$\pm\sigma$ plots
show fault samples carrying elevated spectral power at specific
mid-to-high frequency bins.

%% --- Figure: SpecRas diagnostic outputs ---
\begin{figure}[H]
  \centering
  \subfigure[t-SNE: normal (blue, $n=246$) and fault (orange, $n=154$)
    occupy disjoint clusters with non-overlapping confidence ellipses.]{%
    \includegraphics[width=0.30\linewidth]{cls_tsne.png}%
  }\hfill
  \subfigure[ROC curves: AUC\,=\,0.989 for both classes.
    TPR\,$\approx$\,0.97 at FPR\,=\,0.02.]{%
    \includegraphics[width=0.30\linewidth]{cls_roc.png}%
  }\hfill
  \subfigure[Reliability diagram: ECE\,=\,0.128.
    Bars track the diagonal; model outputs are conservative
    (slight under-confidence).]{%
    \includegraphics[width=0.30\linewidth]{cls_cal.png}%
  }\\[4pt]
  \subfigure[PSD mean\,$\pm1\sigma$ per frequency bin.
    Fault (orange) carries elevated power at mid-to-high frequency
    bands; Fault$-$Normal difference (green dashed) marks the
    discriminative Hz regions.]{%
    \includegraphics[width=0.92\linewidth]{cls_psd.png}%
  }
  \caption{%
    \zerone SWA checkpoint (97.00\% accuracy, macro-F1\,=\,0.969;
    234 TN, 12 FP, 0 FN, 154 TP).
    \textbf{(a)} t-SNE feature clusters.
    \textbf{(b)} ROC (AUC\,=\,0.989).
    \textbf{(c)} Calibration (ECE\,=\,0.128).
    \textbf{(d)} PSD statistics: fault samples show elevated spectral
    power in specific bands identified by the difference curve.
  }
  \label{fig:zerone_results}
\end{figure}

\paragraph{Validation performance}
Both \zerone checkpoints achieve 100\% accuracy and F1 on the
validation set ($N=255$).
The gap to test accuracy (97\%) is attributed entirely to cross-unit
covariate shift, not overfitting.

\subsection{Ablation: SWA Contribution}
\label{sec:ablation}

Table~\ref{tab:ablation} isolates SWA by comparing best-val and SWA
checkpoints, which differ only in the weight-averaging step.

\begin{table}[H]
  \centering
  \small
  \caption{Ablation of stochastic weight averaging (SWA).}
  \label{tab:ablation}
  \renewcommand{\arraystretch}{1.2}
  \begin{tabular}{@{}p{4cm}cc@{}}
    \toprule
    Configuration & Acc.$\uparrow$ & Macro-F1$\uparrow$ \\
    \midrule
    \zerone + ResNet18 (best-val ckpt) & 94.25\% & 0.941 \\
    \zerone + ResNet18 (SWA ckpt)      & \textbf{97.00\%} & \textbf{0.969} \\
    \bottomrule
  \end{tabular}
\end{table}

SWA contributes $+2.75$\,pp accuracy and $+0.028$ macro-F1 at zero
additional inference cost.

\subsection{Misclassification and Visualisation Analysis}
\label{sec:viz}

All 12 false positives from the \zerone SWA checkpoint originate from
transformer 134 and cluster in the 16:00--20:30 local-time window
(peak evening load).
During this period, mag\-net\-o\-stric\-tion-in\-duced broadband excitation is
elevated, pushing high-frequency spectral content into the range that
overlaps genuine fault signatures.
The \hczv false positives (7 samples) overlap the same time window and
show elevated reconstruction error concentrated in the 1\,000--2\,000\,Hz
STFT channel---channel-level root-cause evidence consistent with the
\zerone Grad-CAM attribution maps.

\paragraph{Spectral fingerprint}
Fault samples carry a distinctive PSD pattern: a dominant peak near
$2f_0 = 100$\,Hz and secondary clusters at 300\,Hz and 500\,Hz
(odd harmonics of winding resonance under magnetostriction).
Grad-CAM activations on \texttt{layer4} consistently highlight the
corresponding row bands in the PSD panel (Fig.~\ref{fig:pipeline}).
Fig.~\ref{fig:fp_analysis}(a) confirms this pattern in the PSD
waterfall: the per-sample Fault$-$Normal difference is broadly
positive across 50--300\,Hz, with localised peaks at the harmonic
frequencies.
Fig.~\ref{fig:fp_analysis}(b) dissects a representative \hczv false
positive whose score (1.90) falls just below the threshold (2.00):
reconstruction residuals are largest in the CWT channel at the 100\,Hz
row, consistent with a peak-load magnetostriction harmonic rather than
a genuine fault signature.

%% --- Figure: spectral fingerprint and false-positive analysis ---
\begin{figure}[H]
  \centering
  \subfigure[PSD waterfall (test set).
    Top/middle: per-sample normalised PSD for normal/fault.
    Bottom: Fault$-$Normal difference; red bands mark elevated
    power in the 50--300\,Hz harmonic region.]{%
    \includegraphics[width=0.47\linewidth]{cls_wfall.png}%
  }\hfill
  \subfigure[\hczv false positive (score\,=\,1.90, $\tau=2.00$).
    Columns: input / reconstruction / residual.
    Rows: CWT / \zerone raster / waveform context.
    Residuals concentrate at the 100\,Hz CWT row
    (peak-load magnetostriction, not a structural fault).]{%
    \includegraphics[width=0.47\linewidth]{vae_fp.png}%
  }
  \caption{%
    Spectral fingerprint and false-positive analysis.
    \textbf{(a)} PSD waterfall: Fault$-$Normal difference highlights
    50--300\,Hz harmonic bands; Grad-CAM hot rows in the PSD panel
    correspond to the same frequencies.
    \textbf{(b)} VAE false-positive residuals localise to the 100\,Hz
    CWT row---peak-load magnetostriction excitation, not a fault.
  }
  \label{fig:fp_analysis}
\end{figure}

\paragraph{Training dynamics}
The ReduceLROnPlateau schedule brings the learning rate from
$8\times10^{-6}$ to $8\times10^{-8}$ over 30 epochs.
Validation accuracy and macro-F1 both reach 1.000 by epoch 2 and
remain there, reflecting clean spectral separation between training
normal and fault units.
The 3\,pp drop from validation (100\%) to test (97\%) comes from
cross-unit covariate shift---transformer 134 was recorded under a
different time-of-day distribution, not a harder fault signature.
Fig.~\ref{fig:training_dynamics} shows the training and validation
loss curves together with the validation accuracy and macro-F1 tracks.

%% --- Figure: training dynamics ---
\begin{figure}[H]
  \centering
  \subfigure[Loss vs.\ epoch.
    Train converges to $\approx$0.14 by epoch\,5; val stabilises
    at $\approx$0.17 (cross-unit shift, not overfitting).
    SWA accumulates from epoch\,15 over the flat region.]{%
    \includegraphics[width=0.46\linewidth]{train_loss.png}%
  }\hfill
  \subfigure[Validation accuracy and macro-F1 vs.\ epoch.
    Both metrics saturate at 1.000 from epoch\,2; the flat plateau
    enables stable SWA accumulation.]{%
    \includegraphics[width=0.46\linewidth]{train_score.png}%
  }
  \caption{%
    \zerone ResNet18 training dynamics (30 epochs).
    \textbf{(a)} Loss: rapid convergence; val plateau enables SWA.
    \textbf{(b)} Val accuracy and F1 saturate at 1.000 by epoch\,2.
  }
  \label{fig:training_dynamics}
\end{figure}

Fig.~\ref{fig:overview} shows the dual-pathway architecture and how
energy-weighted channel fusion feeds each diagnostic branch.
Fig.~\ref{fig:pipeline} diagrams the raster grid: the PSD panel
($\approx$77\% of image height) dominates the layout, and a Grad-CAM
hot region at panel row $r$, column $c$ maps to $1 + 32r + c$\,Hz.
Fig.~\ref{fig:signal_chars} illustrates representative multi-channel
vibration characteristics motivating the energy-weighted fusion.
Fig.~\ref{fig:training_dynamics} shows the \zerone training dynamics:
loss convergence by epoch 5 and validation-metric saturation by
epoch 2, enabling stable SWA accumulation from epoch 15.
Fig.~\ref{fig:zerone_results} collects the \zerone diagnostic outputs:
t-SNE feature clusters, ROC curves (AUC\,=\,0.989), calibration
curve (ECE\,=\,0.128), and per-band PSD statistics.
Fig.~\ref{fig:fp_analysis} presents the PSD waterfall and a
representative \hczv false-positive reconstruction breakdown,
localising the residual to the peak-load 100\,Hz CWT row.

\subsection{Discussion}
\label{sec:discussion}

Three findings stand out.

\textit{Why spectral representation matters.}
A model that sees the full shape of the PSD---1\,000 bins rather than
aggregate scalars---can distinguish a peak-load harmonic shift from a
sustained fault signature, whereas coarser features cannot.
The raster encoding makes this tractable for a CNN: with 77\% of the
image area devoted to PSD content, most receptive fields land on
spectral data, and the row-major wrap preserves frequency ordering so
Grad-CAM attribution reads directly as Hz-resolution spectral bands.

\textit{One-class detection versus supervised classification.}
\hczv outperforms the supervised SWA checkpoint by 1.25\,pp
(98.25\% vs.\ 97.00\%).
The spatial VAE optimises a generative target: reconstruct normal
time-frequency images well enough that anything anomalous stands out
by its reconstruction error.
A generative model trained on normal data alone generalises more
reliably to unseen transformer units than a discriminative classifier
fitted on both classes.
The score histogram in Fig.~\ref{fig:vae_results}(c,d) also
shows why the composite score matters: fault and normal distributions
are clearly separated at $\tau=2.00$, with the normal-class peak at
$\approx$4.0 well above the fault-class peak at $\approx$0.7.
The false-positive reconstruction analysis (Fig.~\ref{fig:fp_analysis}(b))
shows that the 7 misclassified normal samples have residuals concentrated
at the 100\,Hz CWT row---a peak-load harmonic, not a fault signature,
and one that the composite reconstruction-plus-Mahalanobis score keeps
just below threshold.

\textit{Load conditions as a confounding variable.}
All false positives from both pathways fall in the 16:00--20:30
peak-load window.
One direct remedy is to condition predictions on ambient load level
by including the transformer tap-changer position or bus voltage as
a covariate.
Alternatively, a time-of-day calibration that widens the normal-class
boundary during known peak-load periods would address it without
requiring additional sensor data.

\textit{Computational overhead.}
Feature extraction (\zerone) processes approximately 4\,000 samples
per CPU-hour.
ResNet18 training completes in under 30 minutes (30 epochs, batch
size 16, GPU).
\hczv training takes approximately 2 hours (50 epochs, batch size 32,
GPU).
Online inference runs in under 100\,ms per Segment, fast enough for
near-real-time monitoring.


%% ========================================================================
\section{Conclusion}
\label{sec:conclusion}
%% ========================================================================

Three things make vibration-based transformer fault diagnosis genuinely
difficult: the signals change with load, confirmed fault examples are
rare, and engineers will not act on a black-box score alone.
This paper addressed all three simultaneously.

The \zerone supervised pathway encodes a 1\,200-dimensional
physics-motivated spectral feature vector as a raster-stripe image and
trains a ResNet18 classifier.
The SWA checkpoint reaches 97.00\% accuracy and macro-F1 of
0.969 on a cross-unit held-out test set, with zero missed faults.
The $+2.75$\,pp gain from SWA over the best-validation checkpoint
confirms that flat-loss-basin regularisation is worth applying whenever
training data are limited and cross-unit covariate shift is expected.

The \hczv pathway trains a SpatialResNetVAE on normal-only data,
retaining a $64{\times}7{\times}7$ spatial latent map that produces
reconstruction-residual heat maps over the time-frequency plane.
On the same test set it reaches 98.25\% accuracy and macro-F1
of 0.982 without using a single fault label.
In practice, substations can go years without a confirmed fault, so a
label-free detector can be deployed immediately rather than waiting for
labelled failures to accumulate.

All false positives from both pathways cluster in the evening peak-load
window and share elevated 1\,000--2\,000\,Hz STFT content, pointing to
a physical rather than a modelling cause.
Incorporating load metadata---tap-changer position or bus voltage---is
the most direct corrective step.

\paragraph{Limitations and future work}
The current evaluation covers two transformer units at one substation
sharing a common operating schedule.
Transfer to units with different core designs (three-phase shell-type
vs.\ core-type), voltage ratings, or 60\,Hz grids remains an open
question; federated or domain-adaptive transfer learning is the
natural approach.
All fault conditions are grouped into a single class---fine-grained
classification of winding deformation, core looseness, and cooling
faults requires labelled examples per fault type, which \hczv
residual heat maps can help accumulate by pointing engineers to
candidate anomaly events.
Finally, Morlet CWT at 128 scales takes roughly 30\,ms per sample on
a modern CPU, too slow for microcontroller-based sensor nodes;
a learnable filter-bank approximation would address this.


%% ---- End matter ----
%% CRediT, Declaration of competing interest, and Acknowledgements
%% have been moved to titlepage.tex (author-identifying information).

\section*{Funding}
This research did not receive any specific grant from funding agencies
in the public, commercial, or not-for-profit sectors.

\section*{Data availability}
The vibration dataset was collected under a confidentiality agreement
with the grid operator and cannot be made publicly available.

%% Acknowledgements are in titlepage.tex.

\section*{Declaration of generative AI and AI-assisted technologies
  in the manuscript preparation process}
During the preparation of this work the author(s) used ChatGPT
(OpenAI) in order to polish the writing and to check grammar and
spelling. After using this tool/service, the author(s) reviewed and
edited the content as needed and take(s) full responsibility for the
content of the published article.

\bibliography{references}

\onecolumn  %% appendices and end-matter figures in single column
\appendix

%% ========================================================================
\section{Feature Extraction Algorithm}
\label{app:algo}
%% ========================================================================

Algorithm~\ref{alg:zerone_feat} summarises the \zerone feature
extraction pipeline for a single multi-channel sample.

\begin{algorithm}[t]
  \small
  \caption{\zerone Feature Extraction}
  \label{alg:zerone_feat}
  \begin{algorithmic}[1]
    \REQUIRE Multi-channel signals $\mathcal{X}_k$,
             sampling rate $f_s$, feature schema $\mathcal{S}$
    \ENSURE Feature vector $\mathbf{S}_k \in \R^{1200}$
    \STATE Compute channel weights $w_u$ via Eq.~\eqref{eq:energy_weight}
    \FOR{each channel $u = 1, \ldots, U_k$}
      \STATE Remove DC: $\bar{x}^{(u)} \gets x_k^{(u)} - \mu_{x_k^{(u)}}$
      \STATE Extract time features $\bs^{(u)}_\mathrm{time}$ (15 dims)
      \STATE Compute Welch PSD via Eq.~\eqref{eq:welch}
      \STATE Extract STFT means $\bs^{(u)}_\mathrm{stft}$ (127 dims)
      \STATE Extract PSD features $\bs^{(u)}_\mathrm{psd}$ (1\,050 dims)
      \STATE Extract HF ratios $\bs^{(u)}_\mathrm{hf}$ (8 dims)
      \STATE Concatenate: $\bs^{(u)} \gets [\bs^{(u)}_\mathrm{time};\,
             \bs^{(u)}_\mathrm{stft};\, \bs^{(u)}_\mathrm{psd};\,
             \bs^{(u)}_\mathrm{hf}]$
    \ENDFOR
    \STATE Aggregate: $\mathbf{S}_k \gets \sum_u w_u \bs^{(u)}$
       via Eq.~\eqref{eq:aggregation}
    \STATE Normalise via Eq.~\eqref{eq:norm}
    \RETURN $\mathbf{S}_k$
  \end{algorithmic}
\end{algorithm}




\end{document}
